\paragraph{Contribution provenance.} The simultaneous constraints continue \texttt{u/UmbrellaCorp\_HR}'s equations (78)--(80), (88)--(90), (96) and (98) \cite{UmbrellaContribution}. This credit identifies the original arithmetic input; the full shared-variable construction below supplies its complete proof and computed lift fibres.

\section{The colleague's shared-variable CRT construction, equations (78)--(98)}
\label{sec:shared-variable-crt}

The additional colleague attachment relayed by the user continues
equations (71)--(100).  It was read in full.  We retain its original
equation numbers in the attributions below.  The results of this
section prove the complete input and output CRT fibres in
(78)--(80) and (96), and construct all the bounded auxiliary
coordinates needed to make (98) an exact theorem with computable
positive integer returns.  No assertion in the attachment is
accepted merely because it has a theorem-status label.

\subsection{Affine input and prescribed output classes}

\begin{proposition}[Complete fibres for source equations (78)--(80)]
\label{prop:input-crt}
Fix a positive integer \(F\) and finitely many
\(\lambda_i,\beta_i\in\Z,\ m_i\in\Z_{>0}\).
The source choice \(F=p!\prod_\nu|E_\nu(p)|\), with
each \(E_\nu(p)\ne0\), is retained as a particular input.
The full input system is
\[
 X\equiv0\pmod{6F},\qquad
 \lambda_iX\equiv\beta_i\pmod{m_i}.
\]
Put \(g_i=\gcd(\lambda_i,m_i)\).  If some \(g_i\nmid\beta_i\),
the system is empty.  Otherwise let
\[
 n_i=m_i/g_i,\qquad
 a_i\equiv(\lambda_i/g_i)^{-1}(\beta_i/g_i)\pmod{n_i},
 \quad 0\le a_i<n_i,
\]
where \(a_i=0\) when \(n_i=1\).
It is solvable if and only if
\[
 \gcd(n_i,n_j)\mid a_i-a_j\quad\text{for every }i,j,
 \qquad \gcd(6F,n_i)\mid a_i\quad\text{for every }i.
\]
When these tests pass, the merge (80), with its stipulated
zero step for reduced modulus one, constructs a unique
residue \(a\in\{0,\ldots,M-1\}\), where
\(M=\lcm(6F,n_1,\ldots,n_s)\).  All integer solutions are
\(a+M\Z\); the associated homogeneous kernel is \(M\Z\).
Its positive solutions and any bounded intersection are exactly
\[
 \{a+Mt:t\in\Z,\ t\ge\lfloor-a/M\rfloor+1\},
\]
\[
 \{a+Mt:\lceil(U-a)/M\rceil\le t\le
                     \lfloor(V-a)/M\rfloor\}\quad(U\le V).
\]
The inverse lift parameter is \(t=(X-a)/M\).
In particular \(X=a+M>0\) satisfies every original condition.
\end{proposition}
\begin{proof}
The coefficient reduction, including \(\lambda_i=0\) and
\(n_i=1\), was proved in Theorem~\ref{thm:divisor-crt}.
The two-residue merge and its full inverse were proved in
Proposition~\ref{prop:crt}.  Necessity of the pairwise tests
follows by subtracting any two required congruences.

Here is a proof that those pairwise tests suffice for the
successive merges, even when the moduli overlap.
Suppose the previously merged modulus is
\(M_0=\lcm(d_0,\ldots,d_k)\), with merged residue \(a_0\);
include \(d_0=6F\) and its zero residue in this list.
For a new residue \(b\pmod n\), consider any prime \(\ell\).
Choose a previous \(d_i\) whose \(\ell\)-valuation is the
largest of the previous valuations.  The common exponent in
\(\gcd(M_0,n)\) is
\(\min(v_\ell(d_i),v_\ell(n))\).
The pairwise condition says that this power divides the
difference between \(b\) and the residue at \(d_i\).
Since \(a_0\) agrees with that residue modulo \(d_i\),
the same power divides \(b-a_0\).
Doing this at every prime proves
\(\gcd(M_0,n)\mid b-a_0\), so the next merge is allowed.
Induction proves all merges.

For the merge, put \(g=\gcd(M_0,n)\) and
\[
 h=\operatorname{rem}_{n/g}
       \left((M_0/g)^{-1}(b-a_0)/g\right).
\]
Then \(a_1=a_0+M_0h\) has both required residues and
the new modulus is \(M_1=M_0n/g\).
If \(0\le a_0<M_0\), the range \(0\le h<n/g\)
gives \(0\le a_1<M_1\).  Thus the claimed representative
range is retained, including \(h=0\) when \(n/g=1\).
The full kernel and fibres follow from the two-residue
proof and induction.  Finally dividing the inequalities
\(0<a+Mt\) or \(U\le a+Mt\le V\) by positive \(M\)
gives precisely the stated integer parameter bounds.
\end{proof}

\begin{proposition}[Prescribed output compatibility and its exact obstruction]
\label{prop:output-crt}
Retain a positive integer seed from the colleague's (88):
\[
 p=12h+1,\quad \varepsilon\in\{0,1\},\quad
 A,B,C,D_p>0,\quad
 4ABCD_p=A+p^{1-\varepsilon}B+pC.
\]
Put
\[
 K_0=4ABC,\quad
 \Theta=C+(1-\varepsilon)B,\quad
 \Omega=A+\varepsilon B,\quad
 g=\gcd(K_0,\Theta),\quad m=K_0/g.
\]
Here \(\Theta\) is the scalar from source equation (88);
it is distinct from the unary bijection of
Theorem~\ref{thm:unary-box}.
For every integer \(q\),
\[
 D_q=(\Omega+q\Theta)/K_0\in\Z
             \quad\Longleftrightarrow\quad q\equiv p\pmod m .
\]
For any \(\mathcal D>0\), the simultaneous output
\(q\equiv1\pmod{\mathcal D}\), \(D_q\in\Z\)
is possible exactly when
\[
 d:=\gcd(\mathcal D,m)\mid p-1.
\]
When this passes, one complete representative and modulus are
\[
 q_0=1+\mathcal D
  \operatorname{rem}_{m/d}
    \left((\mathcal D/d)^{-1}(p-1)/d\right),\qquad
 L_{\rm out}=\lcm(\mathcal D,m).
\]
The remainder is zero when \(m/d=1\).
Every integer output is \(q_0+L_{\rm out}t\), \(t\in\Z\);
its inverse parameter is \((q-q_0)/L_{\rm out}\).
All positive or bounded outputs are obtained by imposing
the corresponding exact linear bounds on \(t\), as above.

The obstruction has the additional exact form
\[
 \operatorname{den}\left(\frac{A+B+C}{4ABC}\right)
       =\frac{m}{\gcd(m,p-1)}>1,
\]
where \(\operatorname{den}\) means the positive denominator
of the reduced rational number.  Thus \(m\nmid p-1\);
\(\mathcal D=12\) is compatible, and any
\(\mathcal D\) divisible by \(m\) is incompatible.
Equivalently, source equation (96) is exactly
\[
 \min(v_\ell(\mathcal D),v_\ell(m))
                      \le v_\ell(p-1)\quad\text{for all primes }\ell.
\]
\end{proposition}
\begin{proof}
The seed identity is \(K_0D_p=\Omega+p\Theta\), so
\(D_q=D_p+(q-p)\Theta/K_0\).
Write \(\Theta=gw\); then \(\gcd(w,m)=1\).
The extra summand is integral exactly when \(m\mid q-p\).
Applying Proposition~\ref{prop:crt} to the two output
classes proves the compatibility test, displayed lift,
homogeneous kernel \(L_{\rm out}\Z\), and its full fibres.
Unique prime factorization proves the valuation form.

At \(q=1\), the same retained identity gives
\[
 \frac{A+B+C}{4ABC}
        =D_p-\frac{(p-1)w}{m}.
\]
Since \(w,m\) are coprime, its reduced denominator is
\(m/\gcd(m,p-1)\).
Each of \(A,B,C\) is at most \(ABC\), and all are positive;
hence the left side lies in \((0,3/4]\).
It is not an integer, which proves the strict inequality
for the denominator and all stated incompatibilities.
For \(\mathcal D=12\), its gcd with \(m\) divides \(12\),
which divides \(p-1\), proving that compatibility.

These calculations prescribe integer output classes.
They neither replace the condition that an output be prime
nor show that a prime divisor of one chosen cyclotomic
value belongs to the prescribed class. The exact incidence
between these two objects, with all input fibres and original
integer witnesses, is computed in the next theorem.
\end{proof}
\input{publication_input_output_incidence.tex}

\subsection{Canonical bounded auxiliary coordinates}

Fix \(p\) throughout the following construction; polynomial
coefficients may be specified integers depending on that fixed \(p\).
Let the original common variables be
\(x=(x_1,\ldots,x_n)\in\Z^n\), with
\[
 l_v\le x_v\le u_v,\qquad
 B_v=\max(|l_v|,|u_v|).
\]
Retain any specified decidable admissibility condition
\(\mathcal A(x)\) on this box, such as the prime label,
the exact shell range, the code decoder, or the valuation budgets.
Its full truth is always tested on one common integer tuple.
Consider a finite Boolean expression \(\Phi\) built from
polynomial equalities, inequalities, and divisibilities
\(d(x)\mid f(x)\), with \(d(x)>0\) on the retained domain.
Comparisons of two polynomials are expressed by their signed
difference, and a congruence is the corresponding divisibility
of that difference.  Implications and biconditionals are
included as Boolean connectives, rather than imposing
either side independently.

For each polynomial \(f=\sum_\alpha c_\alpha x^\alpha\),
define the explicit nonnegative integer bound
\[
 F_f=\sum_\alpha|c_\alpha|\prod_v B_v^{\alpha_v}.
\]
Then \(|f(x)|\le F_f\) throughout the original box.
A factor \(B_v^0\) is one, including when \(B_v=0\).

\begin{lemma}[Unique bounded sign and quotient encoding]
\label{lem:canonical-auxiliary}
The preceding finite truth system admits an explicit polynomial
residual system on a finite augmented box with these properties:
each admissible original tuple has exactly one auxiliary
extension satisfying the residual equations; that extension
records its exact truth bits; and its coordinates have the
following proved bounds.

For a sign test on \(f\), introduce
\(s_-,s_0,s_+\in\{0,1\}\) and
\(0\le t_-,t_+\le\max(F_f-1,0)\).
Use the residual equations
\begin{align}
 s_-+s_0+s_+-1&=0,\notag\\
 f-s_+(t_++1)+s_-(t_-+1)&=0,\label{eq:sign-encoding}\\
 (1-s_+)t_+&=0,\qquad(1-s_-)t_-=0.\notag
\end{align}
The bits for \(f<0,f=0,f>0\) are respectively
\(s_-,s_0,s_+\).

For \(d(x)\mid f(x)\), let \(G=\max(1,F_d)\ge d(x)\)
be the positive integer bound obtained from the absolute
coefficient sum for \(d\) on the original box.  Introduce
\[
 -F_f\le q\le F_f,\qquad 0\le r,s\le G-1,
\]
and residual equations
\begin{equation}\label{eq:quotient-encoding}
 f-dq-r=0,\qquad d-1-r-s=0.
\end{equation}
Then \(q=\lfloor f/d\rfloor\), \(r=f-dq\) is the
canonical remainder and \(s=d-1-r\).  Apply
\eqref{eq:sign-encoding} to \(r\), with bound \(G-1\);
its zero bit is exactly the divisibility bit.
All original Boolean operations are then evaluated once on
these common bits.  Alternatively, introducing one bit per
Boolean subexpression and equating it to its integer truth
polynomial gives a unique bounded polynomial extension as well.
\end{lemma}
\begin{proof}
The first sign equation and the bit domains make exactly one
of \(s_-,s_0,s_+\) equal one.  If it is \(s_+\), the
other equations give \(f=t_++1>0\), \(t_-=0\), and
\(t_+=f-1\).  If it is \(s_-\), they give
\(f=-(t_-+1)<0\), \(t_+=0\), and \(t_-=-f-1\).
If it is \(s_0\), both slacks are zero and \(f=0\).
Conversely these unique choices satisfy all equations and
bounds for each of the three signs; the cases \(F_f=0,1\)
are included by the displayed maximum.  This proves both
existence and uniqueness without an unspecified sign variable.

For division, Euclidean division by \(d>0\) gives unique
integers \(q,r\) with \(f=dq+r,\ 0\le r<d\).
Because \(|f|\le F_f\) and \(d\ge1\), one has
\(-F_f\le\lfloor f/d\rfloor\le F_f\).
Also \(r,s=d-1-r\) lie in \([0,G-1]\).
Conversely the two residual equations and nonnegative
\(r,s\) imply \(0\le r\le d-1\), hence exactly that
canonical Euclidean division, with unique \(s\).
The sign proof applied to \(r\) proves the stated bit.
The finitely many independent auxiliary blocks can therefore
be attached to the same original tuple, each uniquely.
The truth polynomials proved in the parent text take values
zero or one on input bits; evaluating the expression from
its leaves to its root gives a unique bit at every internal
node.  This proves the optional polynomial extension and
its bit bounds by induction on the expression.
\end{proof}

\subsection{The full bounded shared-variable theorem behind equation (98)}

Let \(z=(x,\text{all auxiliaries})\) denote this complete
augmented tuple, whose proved coordinate intervals are
\(\ell_v\le z_v\le u_v\), \(1\le v\le k\).
Keep the original admissibility condition on \(x\) and
all indicated bit domains.  Put
\(B'_v=\max(|\ell_v|,|u_v|)\).
List all residual polynomials from
Lemma~\ref{lem:canonical-auxiliary} and any required
polynomial equalities as
\[
 f_a(z)=\sum_\alpha c_{a,\alpha}(p)z^\alpha.
\]
Required equalities have their own residuals; an equality
appearing inside a negation or biconditional instead has
its computed sign bit and is not incorrectly required true.
Set
\begin{equation}\label{eq:shared-residual-bound}
 H=\max\left(\{0\}\cup
   \left\{\sum_\alpha|c_{a,\alpha}(p)|
                  \prod_v(B'_v)^{\alpha_v}:a\right\}\right).
\end{equation}
Let \(D_{\rm fixed}\ge1\) be any specified common multiple of
all fixed congruence moduli, as in the colleague's equation (98).
Their least common multiple is an explicit possible choice,
with value one for an empty list.  Retain this chosen value and set
\[
 L=D_{\rm fixed}(H+1)>H.
\]
Variable divisor values are handled by
\eqref{eq:quotient-encoding}; they are not treated as
fixed factors of \(D_{\rm fixed}\).

\begin{theorem}[Constructive bounded common lifts and global truth aggregation]
\label{thm:shared-variable-crt}
Choose any nonempty finite list of positive moduli
\(d_1,\ldots,d_s\) whose least common multiple is \(L\).
Overlapping noncoprime moduli are allowed.
Define the vector residue map
\[
 \pi:\Z^k\longrightarrow\prod_{i=1}^s(\Z/d_i\Z)^k.
\]
Its image is exactly the tuples
\(\boldsymbol r=(r_1,\ldots,r_s)\) whose corresponding
coordinates agree modulo every \(\gcd(d_i,d_j)\).
Its additive kernel is \(L\Z^k\).
For each compatible tuple, scalar CRT in each coordinate
constructs its representative
\(r\in\{0,\ldots,L-1\}^k\), and its full integer fibre is
\[
 \pi^{-1}(\boldsymbol r)=r+L\Z^k.
\]
The complete bounded fibre is explicitly
\begin{equation}\label{eq:shared-lift-fibre}
 \left\{r+Lt:
 \left\lceil\frac{\ell_v-r_v}{L}\right\rceil
 \le t_v\le
 \left\lfloor\frac{u_v-r_v}{L}\right\rfloor
 \ \text{for every }v\right\}.
\end{equation}
Its inverse parameter is \(t_v=(z_v-r_v)/L\).
Admissibility and bit-domain restrictions are imposed
on these actual common lifts, retaining every one that passes.

Enumerate exactly those compatible residue tuples for which
\[
 f_a(r_i)\equiv0\pmod{d_i}
       \quad\text{for all residuals }a\text{ and all }i.
\]
For each such tuple, enumerate
\eqref{eq:shared-lift-fibre}, retain the original admissibility
and bit domains, and evaluate the complete Boolean expression
\(\Phi\) on that common lift's canonical bits.
Projection to \(x\) is a bijection from the retained tuples
with \(\Phi=1\) to the original bounded truth-system
solutions.  Its inverse is the unique auxiliary construction
of Lemma~\ref{lem:canonical-auxiliary}.
Thus this algorithm determines existence, all solutions and
their exact integer lifts, including positive ES denominators
when composed with Theorem~\ref{thm:unary-box}.

In particular define
\[
 n(\boldsymbol r)=
 \sum_{\substack{t\text{ in }\eqref{eq:shared-lift-fibre}}}
 [\mathcal A(x(r+Lt))]
 [\text{all specified bit domains at }r+Lt]\,
 [\Phi(\text{recorded bits at }r+Lt)].
\]
For residue tuples that pass the local residual tests above,
the exact solution count is \(\sum_{\boldsymbol r}n(\boldsymbol r)\).
There is no duplicate projection in this count.
The truth fibres, intersections, unions, multiplicities of union
projections and inclusion--exclusion counts of
Theorem~\ref{thm:unary-truth-fibres} hold here on this same
solution domain, with the same common lifted tuple.
\end{theorem}
\begin{proof}
The pairwise sufficiency proof in
Proposition~\ref{prop:input-crt}, applied in each coordinate,
proves the residue image statement.  A vector maps to zero
precisely when every coordinate is divisible by every \(d_i\),
equivalently by \(L\); this proves the kernel.
The coordinatewise CRT representative and kernel give all
integer fibres and the displayed inverse.  Dividing the
original interval inequalities by positive \(L\) proves
\eqref{eq:shared-lift-fibre}, including empty intervals.
The case \(L=1\) has the unique zero residue and
kernel \(\Z^k\), so this description also includes it.

Integer polynomials preserve congruences, since sums and
products preserve them.  Therefore any common lift of a
tuple passing all local residual tests satisfies
\(L\mid f_a(z)\) for every \(a\).
Throughout the proved augmented box,
\eqref{eq:shared-residual-bound} gives
\(|f_a(z)|\le H<L\).  The only multiple of positive
\(L\) in that interval is zero; hence every residual
vanishes as an integer.  Conversely vanishing integer
residuals pass every local residual test.

For an admissible retained lift, the canonical auxiliary
lemma now proves that its recorded bits are the exact
integer truth values of the original predicates, including
all negative remainders and signs.  Thus the single
evaluation \(\Phi=1\) is precisely the original Boolean
requirement.  Every original solution has its unique
auxiliary extension; reducing that extension modulo the
\(d_i\) places it in one of the enumerated compatible
tuples, and its unique lift parameter places it in
\eqref{eq:shared-lift-fibre}.  None is lost.
Conversely each retained lift projects to a solution.
Two retained lifts cannot project to the same \(x\),
because the auxiliary extension is unique, and one
augmented vector has exactly one residue tuple and one
parameter vector \(t\).  This proves both inverse
compositions, the absence of duplicate projections,
and the count.

All products and sums are finite: there are at most
\(L^k\) compatible residue tuples, and each coordinate
interval in the bounded fibre is finite.
Thus this is a terminating exact construction, even if
its enumeration is large.  Applying the previously proved
set-fibre identities on these retained tuples proves
the final union and intersection assertions.
\end{proof}

Two distinctions in the literal content of (98) are essential
and are now calculated, rather than left as lifting assumptions.
First, \(L>H\) makes residual zero-detection exact but need not
make coordinate reduction injective.  For a zero residual,
\(D_{\rm fixed}=2,\ H=0,\ L=2\) and \(-2\le x\le2\),
the zero residue has precisely the three lifts \(-2,0,2\).
Formula \eqref{eq:shared-lift-fibre} returns all three.
Second, a negated composite divisibility is tested after
global aggregation.  In the original shell
\(p=13,a=7,R=15,u=1\), one has \(4u+1=5\):
the congruence passes modulo five and fails modulo three.
Thus \(\neg(15\mid4u+1)\) is true, but
\(\neg(5\mid4u+1)\land\neg(3\mid4u+1)\) is false.
A canonical remainder bit, or the exact global truth
test on the same lift, preserves this difference.

To apply the theorem directly to the complete ES system,
take as original finite domain the unary codes
\(0\le c<N_{\rm code}\) with their exact decoded budget
conditions, or the disjoint tagged valuation boxes of
Theorem~\ref{thm:unary-box}.  Retain every component label,
code, exponent and ordering.  Theorem~\ref{thm:unary-truth-fibres}
evaluates arbitrary additional specified predicates there.
For polynomial residual evaluation, finite decoded
coordinates may instead be included as original bounded
variables with the linear code identity
\eqref{eq:unary-code}, their rectangle bounds, and the
proved integer admissibility tests.
The division and sign construction above supplies every
auxiliary bound required in (98).
The unique code and divisor inverse then gives the original
positive integer denominators, with the exact orientation
in \eqref{eq:unary-integer-lift}.
This proves the requested composition while retaining the
full \(Q_p\) range; \(A\le2\) remains a separate optional
predicate on this same domain.
\begin{corollary}[The actual full primewise system as one polynomial truth system]
\label{cor:actual-prime-truth-system}
For the original prime \(p=12h+1\), use the single rectangle
\[
 0\le j<6h,\quad 1\le A,B\le M_0=9h,\quad
 \varepsilon\in\{0,1\}.
\]
Set, without changing any original shell variable,
\[
 a=3h+j+1,\quad R=4j+3,\quad S=pa,\quad
 T=A+\bigl(p+(1-p)\varepsilon\bigr)B.
\]
The following finite Boolean expression consists solely of
divisibility atoms on that same tuple:
\begin{equation}\label{eq:actual-prime-boolean}
 \Phi_{\rm ES}=
 [AB\mid a]\ \land\
 \bigwedge_{d=2}^{M_0}
       \neg\bigl([d\mid A]\land[d\mid B]\bigr)
       \ \land\ [R\mid T].
\end{equation}
Its successful tuples are exactly the full unary hit set
\(\mathcal C_p^+\) under \eqref{eq:unary-code}.
All divisors in these tests are positive throughout the rectangle.
The canonical quotient construction in
Lemma~\ref{lem:canonical-auxiliary}, applied to \(a/(AB)\)
and \(T/R\), gives the exact variables \(D,C\).
On the truth set they satisfy the explicit bounds
\[
 1\le D\le M_0,\qquad
 1\le C\le C_{\max}:=\left\lfloor\frac{(p+1)M_0}{3}\right\rfloor .
\]
Consequently the original positive integer lifts satisfy
\[
 a\le M_0,\quad u=B^2D\le a^2\le M_0^2,\qquad
 y,z\le pM_0^2C_{\max}.
\]
Appending any specified finite collection of equalities,
positive-divisor congruences, inequalities, implications or
biconditionals on these reconstructed coordinates therefore
gives a fully bounded shared-variable system to which
Theorem~\ref{thm:shared-variable-crt} applies with all
auxiliary bounds constructed above.
The full tagged count and \(2Q_p\) are the respective
unweighted and weighted sums already proved in
Theorem~\ref{thm:unary-truth-fibres}.
\end{corollary}
\begin{proof}
For \(\varepsilon=0,1\), the polynomial
\(p+(1-p)\varepsilon\) equals \(p^{1-\varepsilon}\)
exactly.  Thus \(T\) is the original positive gate numerator.
Also \(a>0,AB>0,R\ge3\), and each fixed \(d\ge2\) is
positive.  If \(\gcd(A,B)>1\), that gcd is an integer in
\([2,M_0]\) dividing both coordinates, so the corresponding
conjunct in \eqref{eq:actual-prime-boolean} fails.
Conversely a failing conjunct supplies a common divisor
greater than one.  The middle conjunction is therefore
exactly \(\gcd(A,B)=1\), proving that the whole expression
is precisely \(\chi_p=1\).

At a hit the two canonical remainders are zero, so their
canonical quotients are \(D=a/(AB)\) and \(C=T/R\).
They are positive because their numerators are positive.
The rectangle gives \(D\le a\le M_0\) and
\(T\le A+pB\le(p+1)M_0\); division by \(R\ge3\)
proves the \(C\) bound.  Since \(a=ABD\),
\[
 u=B^2D\mid A^2B^2D^2=a^2
\]
with exact positive quotient \(A^2D\), proving the
retained divisor budget and its numerical bound.
The reconstruction \eqref{eq:unary-integer-lift},
together with \(p^\varepsilon\le p\) and the established
bounds on \(A,B,C,D\), proves the bounds for both
\(y=p^\varepsilon ACD\) and \(z=pBCD\).
These formulas have already been proved to give the
original ES identity with its exact ordered inverse.

For an appended predicate use these same canonical \(C,D\)
and reconstructed integer coordinates.  Every polynomial
in them has a computable coefficient bound from the
displayed finite intervals; the sign and division lemma
then supplies bounded unique auxiliaries for that predicate.
Equivalently, substitute the coordinate polynomials
\(y=(1+(p-1)\varepsilon)ACD\) and \(z=pBCD\),
retaining their original values and their proved bounds.
The code identity gives a unique code and inverse, so no
additional tuple is introduced by using either presentation.
Applying the shared-variable theorem proves all lifts and
truth fibres; the previously proved multiplicities give
the final count statements.
\end{proof}
